% Appendix: Labb 3, from labs/lab3/README.md: the goals, how the lab is run, the program template,
% the table of parts, the grading and what to do when something does not work. Each of the five
% guides is a section after these, in its own file: first_steps.tex (a_first_steps.md),
% arithmetic.tex (b_arithmetic_and_logic.md), program_flow.tex (c_program_flow.md),
% io_subroutines.tex (d_io_and_subroutines.md) and final_task.tex (e_final_task.md). The
% solutions in labs/lab3/solutions/ are deliberately not in the book.
%
% Deliberate differences: the README shows the template without its header comment; the book
% prints the whole file, labs/lab3/code/template.asm, so the two cannot drift apart. The table of
% parts gives each part's chapter and section instead of its lecture and guide file, and the
% lecture links are chapter and section references.

\chapter{Labb 3: Mikrodatorn}
\label{app:lab3}

\section{Mål}
\label{lab3:sec:goals}
I Labb~3 programmerar du mikrodatorn ATmega328P i assembler. Du börjar med att ladda ett tal i ett
register och slutar med ett trafikljus som körs på ett riktigt Arduino Uno-kort. På vägen använder
du allt kursens mikrodatordel tar upp: talformat, aritmetik och flaggor, logiska operationer, skift,
villkorliga hopp, loopar, tidsfördröjningar, in- och utportar, subrutiner och stacken.

Efter laborationen ska du kunna:
\begin{itemize}
  \item skapa, bygga, rätta och simulera ett assemblerprogram i Microchip Studio;
  \item förutsäga vad ett kort program gör med register, flaggor och portar, och kontrollera det i
    simulatorn;
  \item skriva program med beslut, loopar och tidsfördröjningar utifrån en flödesplan;
  \item programmera mikrodatorn att styra ett styrobjekt: lysdioder som tänds och släcks, och en
    knapp som läses av.
\end{itemize}

\section{Upplägg}
\label{lab3:sec:setup}
Labb~3 löper över de fem mikrodatorpassen, som hör till kapitel~\ref{c7:ch}--\ref{c11:ch}. Varje
pass börjar med en genomgång, och resten av passet är laborationstid. Laborationen består av många
små delar som görs \textbf{i ordning, så långt ni hinner}. Delarna bygger på varandra, så hoppa
inte över någon.
\begin{itemize}
  \item \textbf{Arbeta i par}, vid en dator med Microchip Studio installerat
    (bilaga~\ref{app:studio}). Byt den som skriver vid varje ny del.
  \item \textbf{Ett projekt per del.} Döp projektet efter delen, till exempel \code{lab3_06} för
    del~6, så att ni kan öppna det igen vid redovisningen.
  \item \textbf{Förutsäg innan ni simulerar.} Varje del ber er skriva ned vad ni väntar er innan ni
    stegar. Gör det på riktigt, på papper eller i en kommentar i programmet. Det är förutsägelsen
    som visar om ni har förstått, och en avvikelse är det snabbaste sättet att hitta ett fel
    (\secref{c7:c3}).
  \item \textbf{Redovisa} för läraren vid varje ställe som är markerat
    \kindtag[fillaccent]{Redovisa}. Visa programmet i simulatorn och förklara det; båda i paret ska
    kunna svara på frågor.
\end{itemize}

Nästan allt görs i simulatorn. Arduino-kortet behövs först i slutuppgiften, i
kapitel~\ref{c11:ch}.

\section{Programmallen}
\label{lab3:template}
Varje program utgår från mallen, \file{labs/lab3/code/template.asm}. Kopiera in den i
\code{main.asm} när projektet är skapat, och skriv programmet mellan \code{main:} och \code{end:}.

\asmfile{labs/lab3/code/template.asm}

Vad varje rad gör förklaras i \secref{c7:b6}. Instruktionerna finns samlade på referensbladet i
bilaga~\ref{app:avr}.

\section{Delarna}
\label{lab3:sec:parts}
Numreringen följer den laborationssamling kursen bygger på. Del 13 ingår inte, så numreringen
hoppar från 12 till 14.

\begin{booktable}{@{}p{2.4em}Lp{4em}p{4em}@{}}
  \thead{Del} & \thead{Innehåll} & \thead{Kapitel} & \thead{Avsnitt} \\
  \midrule
  1.1 & Microchip Studio: skapa, skriva, spara, assemblera, simulera & \ref{c7:ch} &
    \ref{lab3:sec:a} \\
  1.2 & Öppna, ändra, rätta, assemblera, simulera & \ref{c7:ch} & \ref{lab3:sec:a} \\
  2 & Decimalt, binärt eller hexadecimalt & \ref{c7:ch} & \ref{lab3:sec:a} \\
  3 & Addition & \ref{c8:ch} & \ref{lab3:sec:b} \\
  4 & Subtraktion & \ref{c8:ch} & \ref{lab3:sec:b} \\
  5 & Addition och subtraktion decimalt, binärt och hexadecimalt & \ref{c8:ch} &
    \ref{lab3:sec:b} \\
  6 & Utporten PORTB, inc, I/O-simulering och .def & \ref{c8:ch} & \ref{lab3:sec:b} \\
  7 & Logiska operationer & \ref{c8:ch} & \ref{lab3:sec:b} \\
  8 & Rotationer och skift & \ref{c8:ch} & \ref{lab3:sec:b} \\
  9 & Aritmetisk rundgång & \ref{c8:ch} & \ref{lab3:sec:b} \\
  10 & Flaggorna i SREG & \ref{c8:ch} & \ref{lab3:sec:b} \\
  11 & Vitsen med flaggor: jämförelse med cp & \ref{c9:ch} & \ref{lab3:sec:c} \\
  12 & Flödesplan och villkorliga hopp & \ref{c9:ch} & \ref{lab3:sec:c} \\
  14 & Loop med visst antal varv & \ref{c9:ch} & \ref{lab3:sec:c} \\
  15 & Kort tidsfördröjning & \ref{c9:ch} & \ref{lab3:sec:c} \\
  16 & Längre tidsfördröjning & \ref{c9:ch} & \ref{lab3:sec:c} \\
  17 & Flervalssituationer och omarbetning av flödesplaner & \ref{c9:ch} & \ref{lab3:sec:c} \\
  18 & ASCII-koder för hexadecimala siffertecken & \ref{c10:ch} & \ref{lab3:sec:d} \\
  19 & 16-bitarsoperationer & \ref{c10:ch} & \ref{lab3:sec:d} \\
  20 & Lång tidsfördröjning & \ref{c10:ch} & \ref{lab3:sec:d} \\
  21 & Inporten PORTD (PIND) & \ref{c10:ch} & \ref{lab3:sec:d} \\
  22 & Subrutiner & \ref{c10:ch} & \ref{lab3:sec:d} \\
  23 & Använda stacken som lagringsplats & \ref{c10:ch} & \ref{lab3:sec:d} \\
  A & Slutuppgift del A: trafikljuset & \ref{c11:ch} & \ref{lab3:sec:e} \\
  B & Slutuppgift del B: övergångsställe & \ref{c11:ch} & \ref{lab3:sec:e} \\
\end{booktable}

Varje pass hör ihop med en del av laborationen, men ni behöver inte bli klara med alla delar på
samma pass. Fortsätt där ni slutade nästa gång.

\section{Redovisning och bedömning}
\label{lab3:sec:grading}

\begin{booktable}{@{}p{5.4em}LL@{}}
  \thead{Laboration} & \thead{Krav för G} & \thead{Krav för VG} \\
  \midrule
  Labb 3 & Del 1.1--12, 14--16, 21 och slutuppgiften del A & Samtliga delar 1.1--23 och
    slutuppgiften del B \\
\end{booktable}

En del är godkänd när programmet är redovisat för läraren i simulatorn, eller på kortet i
slutuppgiften, och båda i paret kan förklara det. Skriv upp vilka delar som är redovisade, och
spara projekten tills kursen är slut.

\section{När något inte fungerar}
\label{lab3:sec:trouble}
\begin{itemize}
  \item \textbf{Bygg om och läs det första felet} i Error List (\secref{c7:c4}).
  \item \textbf{Stega fram till felet.} Leta upp den första instruktion där registren inte blir det
    ni förutsade. Felet sitter där eller strax före.
  \item \textbf{Ändra en sak i taget}, och bygg om efter varje ändring.
  \item \textbf{Fråga.} Läraren hjälper hellre till tidigt än ser er fastna i en halvtimme.
\end{itemize}

Vanliga fel samlas i tabellen i \secref{studio:6}.

% Labb 3, parts 1.1, 1.2 and 2, from labs/lab3/a_first_steps.md. Each part is a \labtask, and its
% Mål, Bakgrund, Programmet, Uppgifter and Frågor are \task headings. The program given in part
% 1.1 and the four faulty lines in part 1.2 are the handout's, and are printed as the guide
% prints them.

\section{Del 1.1, 1.2 och 2: Första stegen}
\label{lab3:sec:a}
Delarna hör till kapitel~\ref{c7:ch}. Läs \secref{c7:app:b} och \secref{c7:app:c} innan ni börjar,
och ha bilaga~\ref{app:studio} om Microchip Studio öppen bredvid.

\labtask{Del 1.1}{Microchip Studio: skapa, skriva, spara, assemblera, simulera}%
\phantomsection\label{lab3:d1-1}
\task{Mål}
Skapa ett assemblerprojekt, skriva ett program, bygga det och stega igenom det i simulatorn.

\task{Bakgrund}
Arbetsgången, skriv, bygg, förutsäg, stega och jämför, beskrivs i \secref{c7:c1}. Instruktionerna
i programmet står i \secref{c7:b5}.

\task{Programmet}
\begin{asmcode}
.include "m328Pdef.inc"

.org 0x0000
    rjmp main

main:
    ldi r16, 10                     ; r16 = 10
    ldi r17, 20                     ; r17 = 20
    mov r18, r16                    ; r18 = r16
    inc r18                         ; r18 = r18 + 1
    inc r18                         ; r18 = r18 + 1
    dec r17                         ; r17 = r17 - 1

end:
    rjmp end
\end{asmcode}

\task{Uppgifter}
\begin{enumerate}
  \item Skapa ett projekt med namnet \code{lab3_01_1} enligt \secref{studio:2}.
  \item Ersätt innehållet i \code{main.asm} med programmet ovan. \textbf{Skriv av det} i stället
    för att klistra in: det är så ni lär er hur en rad är uppbyggd. Kommentarerna kan ni hoppa
    över.
  \item Spara med \textbf{Ctrl+S}.
  \item Bygg med \textbf{F7}. Läs utskriften i fönstret Output: står det \code{0 errors}? Om inte,
    rätta raden Error List pekar på och bygg om.
  \item Välj simulatorn som verktyg (\secref{studio:4-1}) och starta med \textbf{Debug
    $\rightarrow$ Start Debugging and Break} (\textbf{Alt+F5}). Öppna fönstret \textbf{Processor
    Status}.
  \item \textbf{Innan ni stegar:} fyll i tabellen nedan med registrens värden efter varje
    instruktion, decimalt och hexadecimalt.

    \begin{filltable}{@{}p{8.4em}YYY@{}}
      \thead{Efter} & \thead{\code{r16}} & \thead{\code{r17}} & \thead{\code{r18}} \\
      \midrule
      \code{ldi r16, 10} & & & \fillrule
      \code{ldi r17, 20} & & & \fillrule
      \code{mov r18, r16} & & & \fillrule
      \code{inc r18} & & & \fillrule
      \code{inc r18} & & & \fillrule
      \code{dec r17} & & & \\
    \end{filltable}
  \item Stega med \textbf{F11}, en instruktion i taget, och jämför Processor Status med tabellen
    efter varje steg. Lägg märke till att registret som just ändrades visas i rött.
  \item Läs av \textbf{Program Counter} vid varje steg. Med hur mycket ökar den för varje
    instruktion?
  \item När den gula pilen står på \code{rjmp end}: vad visar \textbf{Cycle Counter}? Räkna ut
    värdet för hand först, med hjälp av \secref{c7:c5}.
\end{enumerate}

\task{Frågor}
\begin{itemize}
  \item Den gula pilen står på \code{inc r18}. Har \code{inc r18} körts ännu?
  \item Ändras \code{r16} av instruktionen \code{mov r18, r16}?
\end{itemize}

\redovisa{tabellen och simuleringen för läraren.}

\labtask{Del 1.2}{Öppna, ändra, rätta, assemblera, simulera}\phantomsection\label{lab3:d1-2}
\task{Mål}
Öppna ett befintligt projekt, ändra programmet, och rätta fel med hjälp av felmeddelandena.

\task{Bakgrund}
De vanligaste felmeddelandena, och hur de rättas, står i \secref{c7:c4}.

\task{Uppgifter}
\begin{enumerate}
  \item Stäng Microchip Studio och starta det igen. Öppna projektet från del 1.1, antingen från
    listan över senaste projekt på startsidan eller med \textbf{File $\rightarrow$ Open
    $\rightarrow$ Project/Solution...}
  \item Ändra programmet:
    \begin{itemize}
      \item \code{r16} ska laddas med 25 i stället för 10;
      \item lägg till raden \code{mov r19, r17} efter \code{dec r17}.
    \end{itemize}
  \item \textbf{Förutsäg} värdena i \code{r16}--\code{r19} när programmet når \code{end}. Bygg,
    simulera och kontrollera.
  \item Lägg nu till de här fyra raderna direkt efter \code{mov r19, r17}. De innehåller fyra fel
    med flit:
\begin{asmcode}
    ldi r16 7
    ldi r5, 3
    inc r19, 2
    rjmp ennd
\end{asmcode}
    Meningen med raderna är: ladda 7 i \code{r16}, ladda 3 i något register, öka \code{r19} med
    två, och hoppa till \code{end}.
  \item Bygg. Läs det första felet i \textbf{Error List}, dubbelklicka på det, rätta raden och bygg
    om. Fortsätt tills programmet byggs utan fel. Skriv för varje fel ned vad meddelandet sa och
    hur ni rättade det.
  \item \textbf{Förutsäg} värdena i \code{r16}--\code{r20} när programmet når \code{end}. Simulera
    och kontrollera.
\end{enumerate}

\task{Frågor}
\begin{itemize}
  \item Hur många gånger behövde ni bygga om innan alla fel var borta? Fick ni alla fyra felen på
    en gång?
  \item Varför kan \code{inc} inte öka med två på en gång?
\end{itemize}

\redovisa{det rättade programmet och de fyra felmeddelandena.}

\labtask{Del 2}{Decimalt, binärt eller hexadecimalt}\phantomsection\label{lab3:d2}
\task{Mål}
Skriva tal i ett program decimalt, binärt och hexadecimalt, och läsa dem i simulatorn.

\task{Bakgrund}
Talformaten står i \secref{c7:b4} och omvandlingen mellan talsystemen i kapitel~\ref{c1:ch}.

\task{Uppgifter}
\begin{enumerate}
  \item Skriv ett program som laddar talet 200 i \code{r16} decimalt, i \code{r17} binärt och i
    \code{r18} hexadecimalt. Räkna ut det binära och det hexadecimala skrivsättet för hand.
    \textbf{Förutsäg} hur Processor Status visar de tre registren. Simulera och kontrollera.
  \item Fyll i de tomma rutorna i tabellen \textbf{för hand}:

    \begin{filltable}{@{}p{5em}YYY@{}}
      \thead{Register} & \thead{Decimalt} & \thead{Binärt} & \thead{Hexadecimalt} \\
      \midrule
      \code{r19} & 9 & & \fillrule
      \code{r20} & & \code{0b00001111} & \fillrule
      \code{r21} & & & \code{0x10} \fillrule
      \code{r22} & 100 & & \fillrule
      \code{r23} & & \code{0b01111111} & \fillrule
      \code{r24} & & & \code{0x80} \fillrule
      \code{r25} & 255 & & \\
    \end{filltable}
  \item Lägg till en \code{ldi} per rad i programmet, som laddar registret med talet \textbf{i det
    skrivsätt tabellen ger det}: \code{ldi r19, 9}, \code{ldi r20, 0b00001111}, och så vidare.
    Simulera och kontrollera att Processor Status visar det hexadecimala värde ni räknade ut på
    varje rad.
  \item Lägg till \code{ldi r26, '0'} och \code{ldi r27, 'A'}. Vilka värden får registren? Ett
    tecken lagras som ett tal, dess ASCII-kod; det återkommer i del~18.
  \item Välj för varje rad i tabellen vilket skrivsätt som passar bäst om talet är
    \begin{itemize}
      \item ett antal varv i en loop;
      \item ett mönster av lysdioder som ska lysa.
    \end{itemize}
\end{enumerate}

\task{Frågor}
\begin{itemize}
  \item Blir maskinkoden olika för \code{ldi r16, 200} och \code{ldi r16, 0xC8}? Titta i listfilen
    om ni är osäkra (\secref{c7:b7}).
\end{itemize}

\redovisa{tabellen och simuleringen.}

% Labb 3, parts 3 to 10, from labs/lab3/b_arithmetic_and_logic.md. Each part is a \labtask, and
% its Mål, Bakgrund, Programmet, Uppgifter and Frågor are \task headings. The program given in
% part 6 is the handout's.

\section{Del 3--10: Aritmetik, logik och flaggor}
\label{lab3:sec:b}
Delarna hör till kapitel~\ref{c8:ch}. Varje del är ett eget projekt som utgår från mallen
(\secref{lab3:template}). Ha referensbladet i bilaga~\ref{app:avr} till hands.

\textbf{Förutsäg alltid först.} Varje del har uppgifter som börjar med \emph{Förutsäg}. Skriv svaret
på papper eller som en kommentar i programmet innan ni stegar, och jämför sedan.

\labtask{Del 3}{Addition}\phantomsection\label{lab3:d3}
\task{Mål}
Addera med \code{add} och \code{inc}, och se vad som händer när summan inte får plats i en byte.

\task{Bakgrund}
\Secref{c8:a1} och~\ref{c8:a2}.

\task{Uppgifter}
\begin{enumerate}
  \item Skriv ett program som laddar 25 i \code{r16} och 17 i \code{r17}, och adderar dem med
    \code{add r16, r17}. \textbf{Förutsäg} \code{r16} decimalt, hexadecimalt och binärt. Simulera
    och kontrollera.
  \item Gör samma addition för hand på papper, binärt, kolumn för kolumn med minnessiffror.
  \item Lägg till en addition av två hexadecimala tal: \code{0x3A} i \code{r18} plus \code{0x15} i
    \code{r19}. \textbf{Förutsäg} \code{r18}, både hexadecimalt och decimalt.
  \item Lägg till instruktioner som räknar ut 10 + 20 + 30 i \code{r20}. Hur många \code{add}
    behövs?
  \item Lägg till en addition av 200 i \code{r23} och 100 i \code{r24}. \textbf{Förutsäg}
    \code{r23}. Titta sedan på flaggan \code{C} i Processor Status efter additionen. Vad betyder den
    här?
\end{enumerate}

\task{Frågor}
\begin{itemize}
  \item Vilket av de två registren i \code{add r16, r17} ändras?
  \item Summan 200 + 100 = 300 får inte plats. Var finns \enquote{resten} av svaret efter
    additionen?
\end{itemize}

\redovisa{programmet, uträkningen från uppgift~2 och svaret på frågorna.}

\labtask{Del 4}{Subtraktion}\phantomsection\label{lab3:d4}
\task{Mål}
Subtrahera med \code{sub}, \code{subi} och \code{dec}, och addera en konstant utan en
\code{addi}.

\task{Bakgrund}
\Secref{c8:a4}.

\task{Uppgifter}
\begin{enumerate}
  \item Räkna ut 50 -- 20 med \code{sub}, med 50 i \code{r16} och 20 i \code{r17}.
    \textbf{Förutsäg} resultatet.
  \item Ladda \code{0x20} i \code{r18} och dra ifrån \code{0x05} med \code{subi}. \textbf{Förutsäg}
    \code{r18} hexadecimalt och decimalt.
  \item Ladda 3 i \code{r19} och kör \code{dec r19} tre gånger. \textbf{Förutsäg} \code{r19} och
    flaggan \code{Z} efter varje \code{dec}. Sätt en etikett efter den tredje, så kan ni köra dit
    med \textbf{Run To Cursor}.
  \item Det finns ingen \code{addi}. Ladda 100 i \code{r20} och addera 28 med en \code{subi}.
    Vilken konstant ska \code{subi} ha? \textbf{Förutsäg} \code{r20}.
  \item Räkna ut 20 -- 50 med \code{sub}, med 20 i \code{r21} och 50 i \code{r22}.
    \textbf{Förutsäg} \code{r21} och \code{C}.
\end{enumerate}

\task{Frågor}
\begin{itemize}
  \item Vad betyder \code{C = 1} efter en subtraktion?
  \item \code{r21} blev 226 i uppgift~5. Vilket tal är det om man läser det med tecken? Jämför med
    \secref{c8:a5}.
\end{itemize}

\redovisa{programmet och förutsägelserna.}

\labtask{Del 5}{Addition och subtraktion decimalt, binärt och hexadecimalt}%
\phantomsection\label{lab3:d5}
\task{Mål}
Räkna med tal som är skrivna i olika talsystem i samma uttryck.

\task{Bakgrund}
\Secref{c7:b4} och kapitel~\ref{c1:ch}.

\task{Uppgifter}
Skriv ett program som räknar ut de tre uttrycken, med talen skrivna \textbf{precis som i
uttrycket}. \textbf{Förutsäg} först varje resultat decimalt, hexadecimalt och binärt, i en tabell.
\begin{enumerate}
  \item \code{r20} = \code{0x3C} + \code{0b00001010} -- 25
  \item \code{r21} = \code{0b11110000} -- \code{0x0F} + 1, där den sista additionen görs med
    \code{inc}
  \item \code{r22} = \code{0b10101010} + \code{0x55}
\end{enumerate}

\begin{filltable}{@{}p{5em}YYY@{}}
  \thead{Register} & \thead{Decimalt} & \thead{Hexadecimalt} & \thead{Binärt} \\
  \midrule
  \code{r20} & & & \fillrule
  \code{r21} & & & \fillrule
  \code{r22} & & & \\
\end{filltable}

Simulera och jämför med tabellen.

\task{Frågor}
\begin{itemize}
  \item Uttryck~3 ger ett särskilt bitmönster. Vilket, och varför blev det så? Titta på de två
    talen binärt.
\end{itemize}

\redovisa{tabellen och simuleringen.}

\labtask{Del 6}{Utporten PORTB, inc, I/O-simulering och .def}\phantomsection\label{lab3:d6}
\task{Mål}
Göra port B till utport, räkna på den, och följa bitarna i simulatorns I/O-fönster.

\task{Bakgrund}
\Secref{c8:app:c}, särskilt \ref{c8:c3}, \ref{c8:c6} och~\ref{c8:c7}.

\task{Programmet}
\begin{asmcode}
.include "m328Pdef.inc"

.def counter = r16                  ; The value shown on port B.
.def temp = r17                     ; A scratch register.

.org 0x0000
    rjmp main

main:
    ser temp                        ; temp = 0xFF
    out DDRB, temp                  ; All eight pins of port B are outputs.
    clr counter                     ; Start counting at zero.

loop:
    out PORTB, counter              ; Show the count on port B.
    inc counter                     ; Count up.
    rjmp loop                       ; And again, forever.
\end{asmcode}

\task{Uppgifter}
\begin{enumerate}
  \item Skriv av programmet, bygg och starta simuleringen. Öppna fönstret \textbf{I/O}
    (\textbf{Debug $\rightarrow$ Windows $\rightarrow$ I/O}) och välj port B.
  \item Stega och följ rutorna i \code{DDRB} och \code{PORTB}. När blir alla rutor i \code{DDRB}
    fyllda?
  \item \textbf{Förutsäg} hur rutorna i \code{PORTB} ser ut efter fem varv i loopen. Stega och
    kontrollera.
  \item Varför ändras även \code{PINB}, fast programmet aldrig skriver dit?
  \item Ändra programmet så att det räknar \textbf{nedåt}. \textbf{Förutsäg} vilket värde som
    visas direkt efter 0.
  \item Ändra programmet så att det räknar uppåt i steg om \textbf{två}, med en enda instruktion i
    stället för \code{inc}.
  \item Kör programmet fritt med \textbf{F5}, vänta några sekunder och stoppa med \textbf{Break
    All}. Går det att förutsäga värdet i \code{PORTB}? Hur lång tid tar ett varv i loopen, vid
    16~MHz?
\end{enumerate}

\task{Frågor}
\begin{itemize}
  \item Vad gör \code{.def}, och blir det någon skillnad i maskinkoden?
  \item Hur skulle lysdioderna se ut på ett riktigt kort om programmet kördes fritt?
\end{itemize}

\redovisa{programmet med nedräkning och med steg om två, och svaren på frågorna.}

\labtask{Del 7}{Logiska operationer}\phantomsection\label{lab3:d7}
\task{Mål}
Använda \code{andi}, \code{ori}, \code{eor} och \code{com} med masker för att ändra enskilda
bitar.

\task{Bakgrund}
\Secref{c8:b1} och~\ref{c8:b2}.

\task{Uppgifter}
Gör port B till utport först i programmet, och skriv varje resultat till \code{PORTB} med
\code{out}, så att det syns bit för bit i I/O-fönstret. \textbf{Förutsäg} varje resultat binärt
innan ni stegar.
\begin{enumerate}
  \item Ladda \code{0b10101010} i \code{r16} och behåll bara de fyra nedre bitarna.
  \item Ladda \code{0b00110000} i \code{r17} och ettställ bit~7 och bit~0.
  \item Ladda \code{0b00001111} i \code{r18} och invertera de fyra övre bitarna. Det finns ingen
    \code{eori}: hur gör ni?
  \item Ladda \code{0b00110011} i \code{r19} och invertera alla bitar med en enda instruktion.
  \item Ladda \code{0b01010101} i \code{r20}. Ettställ bit~1, nollställ bit~6 och invertera bit~7,
    med en instruktion per ändring. Vilket värde får \code{r20}, och vilka stift på port B lyser?
\end{enumerate}

\task{Frågor}
\begin{itemize}
  \item Vilken instruktion och vilken sorts mask använder man för att nollställa bitar? För att
    ettställa?
  \item Varför ger \code{com} och \code{eor} med \code{0xFF} samma resultat?
\end{itemize}

\redovisa{programmet och förutsägelserna.}

\labtask{Del 8}{Rotationer och skift}\phantomsection\label{lab3:d8}
\task{Mål}
Skifta och rotera bitar, och använda skift för att multiplicera och dividera med 2.

\task{Bakgrund}
\Secref{c8:b3} och~\ref{c8:b4}.

\task{Uppgifter}
\textbf{Förutsäg} varje resultat, och \code{C}, innan ni stegar.
\begin{enumerate}
  \item Ladda 3 i \code{r16} och skifta vänster två gånger med \code{lsl}. Vilket tal har
    \code{r16} multiplicerats med?
  \item Ladda 200 i \code{r17} och skifta höger två gånger med \code{lsr}.
  \item Ladda \code{0b11110000} i \code{r18}, som är $-16$ med tecken, och kopiera det till
    \code{r19}. Skifta \code{r18} med \code{asr} och \code{r19} med \code{lsr}. Vilket av
    resultaten är $-16 / 2$?
  \item Ladda \code{0b10010011} i \code{r20}, nollställ \code{C} med \code{clc}, och rotera vänster
    två gånger med \code{rol}. \textbf{Förutsäg} \code{r20} och \code{C} efter varje \code{rol}.
  \item Skriv ett rinnande ljus på port B med \code{rol r21}, som i \secref{c8:b4}. Stega igenom
    ett helt varv, och följ \code{PORTB} och \code{C}. Hur många steg går det på ett varv, och
    varför är ett av dem mörkt?
\end{enumerate}

\task{Frågor}
\begin{itemize}
  \item Varför används \code{asr} och inte \code{lsr} för att dividera ett tal med tecken med 2?
  \item \emph{(Fördjupning)} Byt ut \code{rol r21} mot de två instruktionerna \code{lsl r21} och
    \code{adc r21, r25}, där \code{r25} har nollställts med \code{clr} före loopen. Hur många steg
    går ljuset nu på ett varv, och varför?
\end{itemize}

\redovisa{det rinnande ljuset och förutsägelserna.}

\labtask{Del 9}{Aritmetisk rundgång}\phantomsection\label{lab3:d9}
\task{Mål}
Se vad som händer när ett register räknas förbi 255 eller under 0, och vilka flaggor som visar
det.

\task{Bakgrund}
\Secref{c8:a3}.

\task{Uppgifter}
Sätt en etikett efter varje uppgift, så att ni kan köra fram till den och läsa SREG.
\textbf{Förutsäg} registret och flaggorna \code{Z} och \code{C} innan ni kör.
\begin{enumerate}
  \item Nollställ \code{C} med \code{clc}. Ladda 250 i \code{r16} och kör \code{inc r16} sex
    gånger.
  \item Ladda 255 i \code{r17} och 1 i \code{r18}, och addera dem med \code{add}.
  \item Ladda 0 i \code{r19} och kör \code{dec r19}. Titta även på \code{N}.
  \item Ladda 3 i \code{r20} och 5 i \code{r21}, och räkna ut \code{r20 - r21} med \code{sub}.
  \item Ladda 200 i \code{r22} och 100 i \code{r23}, och addera dem.
\end{enumerate}

\task{Frågor}
\begin{itemize}
  \item Efter uppgift~1 är \code{r16} = 0 och \code{Z} = 1. Vad är \code{C}, och varför?
  \item Efter uppgift~3 är \code{C} = 1. Är det \code{dec} som har satt den? Vilken instruktion
    satte den?
  \item Vilket tal ska man lägga till 44 för att få det \enquote{riktiga} svaret på uppgift~5?
\end{itemize}

\redovisa{förutsägelserna och svaren på frågorna.}

\labtask{Del 10}{Flaggorna i SREG}\phantomsection\label{lab3:d10}
\task{Mål}
Förutsäga flaggorna \code{C}, \code{Z}, \code{N}, \code{V}, \code{S} och \code{H} efter olika
operationer, och kontrollera dem.

\task{Bakgrund}
\Secref{c8:a6} och~\ref{c8:a7}.

\task{Uppgifter}
\begin{enumerate}
  \item Fyll i tabellen \textbf{för hand}. Varje rad görs för sig, med \code{add}, \code{sub},
    \code{andi} eller \code{com} som anges.

    \begin{filltable}{@{}cp{11.2em}Y*{6}{C{1.5em}}@{}}
      \thead{Rad} & \thead{Operation} & \thead{Resultat} & \thead{C} & \thead{Z} & \thead{N} &
        \thead{V} & \thead{S} & \thead{H} \\
      \midrule
      a & \code{0x7F + 0x01} & & & & & & & \fillrule
      b & \code{0xFF + 0x01} & & & & & & & \fillrule
      c & \code{0x80 + 0x80} & & & & & & & \fillrule
      d & \code{0x05 - 0x07} & & & & & & & \fillrule
      e & \code{0x40 + 0x40} & & & & & & & \fillrule
      f & \code{0x0F + 0x01} & & & & & & & \fillrule
      g & \code{0x0F} AND \code{0xF0} (\code{andi}) & & & & & & & \fillrule
      h & NOT \code{0x00} (\code{com}) & & & & & & & \\
    \end{filltable}
  \item Skriv ett program med en rad per operation och en etikett efter varje, till exempel
    \code{row_a}, \code{row_b} och så vidare.
  \item Kör fram till varje etikett och jämför SREG med tabellen. Markera varje ruta ni hade fel
    på.
  \item För rad g och h: vilka flaggor ändrades inte av operationen? Jämför med kolumnen
    \emph{Flaggor} på referensbladet i bilaga~\ref{app:avr}.
\end{enumerate}

\task{Frågor}
\begin{itemize}
  \item Rad a och rad e ger båda \code{V = 1}. Vad har de gemensamt?
  \item Rad d ger \code{C = 1} men \code{V = 0}. Är svaret rätt eller fel? Svara både utan och med
    tecken.
\end{itemize}

\redovisa{tabellen, med era förutsägelser och de rätta värdena bredvid varandra.}

% Labb 3, parts 11, 12 and 14 to 17, from labs/lab3/c_program_flow.md. Each part is a \labtask,
% and its Mål, Bakgrund, Specifikation, Programmet, Uppgifter and Frågor are \task headings. The
% program given in part 15 is the handout's.

\section{Del 11, 12 och 14--17: Programflöde, loopar och tidsfördröjningar}
\label{lab3:sec:c}
Delarna hör till kapitel~\ref{c9:ch}. Varje del är ett eget projekt som utgår från mallen
(\secref{lab3:template}).

Från och med nu gäller två vanor till:
\begin{itemize}
  \item \textbf{Rita flödesplanen först}, på papper, innan ni skriver en enda instruktion
    (\secref{c9:a1}).
  \item \textbf{Mät tider med stoppuret}, med Frequency inställd på 16~MHz (\secref{c9:b6}).
\end{itemize}

Del 13 ingår inte i laborationen; numreringen hoppar från 12 till 14.

\labtask{Del 11}{Vitsen med flaggor: jämförelse med cp}\phantomsection\label{lab3:d11}
\task{Mål}
Se vad \code{cp} och \code{cpi} gör med flaggorna, och hur flaggorna säger om två tal är lika,
större eller mindre.

\task{Bakgrund}
\Secref{c9:a3} och~\ref{c9:a6}.

\task{Uppgifter}
Skriv ett program med de fem jämförelserna nedan, med en etikett efter var och en.
\textbf{Förutsäg} \code{Z}, \code{C} och \code{S} för varje jämförelse i en tabell innan ni kör
fram till etiketterna.

\begin{filltable}{@{}p{8.8em}cc*{3}{C{1.8em}}Y@{}}
  \thead{Jämförelse} & \thead{\code{r16}} & \thead{\code{r17}} & \thead{Z} & \thead{C} &
    \thead{S} & \thead{\code{r16} är \dots\ \code{r17}} \\
  \midrule
  1. \code{cp r16, r17} & 50 & 30 & & & & \fillrule
  2. \code{cp r16, r17} & 50 & 50 & & & & \fillrule
  3. \code{cp r16, r17} & 50 & 70 & & & & \fillrule
  4. \code{cpi r16, 50} & 50 & -- & & & & \fillrule
  5. \code{cp r16, r17} & $-5$ & 3 & & & & \\
\end{filltable}

\begin{enumerate}
  \item Fyll i tabellen, och fyll i sista kolumnen med \emph{större än}, \emph{lika med} eller
    \emph{mindre än}.
  \item Kör fram till varje etikett och kontrollera flaggorna.
  \item Kontrollera efter jämförelse 1--3 att \code{r16} fortfarande är 50. Varför ändras det inte,
    fast \code{cp} räknar ut en subtraktion?
  \item Titta särskilt på jämförelse~5. Är $-5$ mindre eller större än 3? Svara både
    \textbf{utan} tecken, när \code{r16} läses som 251, och \textbf{med} tecken. Vilken flagga
    svarar på vilken fråga?
\end{enumerate}

\task{Frågor}
\begin{itemize}
  \item Vilka två flaggor räcker för att avgöra om två tal utan tecken är lika, om det första är
    större, eller om det är mindre?
  \item Varför behövs både \code{brlo} och \code{brlt}?
\end{itemize}

\redovisa{tabellen.}

\labtask{Del 12}{Flödesplan och villkorliga hopp}\phantomsection\label{lab3:d12}
\task{Mål}
Rita en flödesplan med beslut, och skriva det som ett program med jämförelser och villkorliga hopp.

\task{Bakgrund}
\Secref{c9:a4} och~\ref{c9:a5}.

\task{Specifikation}
En termostat styr ett rum. Temperaturen i grader finns i \code{r16}. Port B styr två saker:
\begin{itemize}
  \item \textbf{PB0}, en värmare, ska vara på när temperaturen är \textbf{under 20} grader;
  \item \textbf{PB1}, en fläkt, ska vara på när temperaturen är \textbf{25 grader eller mer};
  \item mellan 20 och 24 grader ska båda vara av.
\end{itemize}

Programmet ska gå i en loop: det läser \code{r16}, sätter utgångarna, och börjar om. Då kan ni
stoppa simuleringen, ändra värdet i \code{r16} i Processor Status, och köra vidare för att se
utgångarna ändras.

\task{Uppgifter}
\begin{enumerate}
  \item Rita flödesplanen. Den har två beslut. Märk varje beslut med Ja och Nej.
  \item Skriv programmet. Ladda startvärdet 18 i \code{r16} \textbf{före} loopen, så att ett värde
    ni skriver in i Processor Status inte skrivs över.
  \item \textbf{Förutsäg} \code{PORTB} för temperaturerna 18, 20, 24, 25 och 30. Testa alla fem i
    simulatorn.
  \item Vilka av värdena är \textbf{gränsfall}, och varför är det de som är viktigast att testa?
\end{enumerate}

\task{Frågor}
\begin{itemize}
  \item Villkoret \enquote{25 grader eller mer} används för fläkten. Hur skriver man \enquote{mer
    än 24} med de hopp som finns?
  \item Vad händer med termostaten vid $-5$ grader? Se \secref{c9:a6}.
\end{itemize}

\redovisa{flödesplanen och programmet, och visa minst tre temperaturer i simulatorn.}

\labtask{Del 14}{Loop med visst antal varv}\phantomsection\label{lab3:d14}
\task{Mål}
Skriva en loop som går ett bestämt antal varv med en räknare, och veta exakt hur många gånger
varje instruktion i den körs.

\task{Bakgrund}
\Secref{c9:b1}.

\task{Specifikation}
Processorn har ingen multiplikation i den här laborationen. Räkna ut $7 \cdot 3$ genom att addera 3
sju gånger, med en loop. Visa antalet varv som hittills gått på port B, så att det syns i
I/O-fönstret.

\task{Uppgifter}
\begin{enumerate}
  \item Rita flödesplanen. Använd ett register för resultatet, ett för talet som adderas, ett för
    varven som är kvar och ett för varven som har gått.
  \item Skriv programmet med \code{.equ} för konstanterna, till exempel \code{.equ LAPS = 7} och
    \code{.equ TERM = 3}.
  \item \textbf{Förutsäg} resultatet och \code{PORTB} när programmet är klart. Stega igenom minst
    två varv och se \code{PORTB} räkna, kör sedan till slutet med \textbf{Run To Cursor}.
  \item Hur många gånger körs \code{brne}, och hur många av dem hoppar den?
  \item Ändra \code{LAPS} till 0. \textbf{Förutsäg} resultatet innan ni kör. Förklara det ni ser.
  \item Räkna ut $12 \cdot 25$ med samma program. Stämmer resultatet? Varför, eller varför inte?
\end{enumerate}

\task{Frågor}
\begin{itemize}
  \item Varför behövs ingen \code{cpi} före \code{brne}?
\end{itemize}

\redovisa{programmet och svaren på uppgift 4--6.}

\labtask{Del 15}{Kort tidsfördröjning}\phantomsection\label{lab3:d15}
\task{Mål}
Göra en tidsfördröjning med en loop, räkna ut dess längd i klockcykler, och kontrollera den med
stoppuret.

\task{Bakgrund}
\Secref{c9:b3} och~\ref{c9:b4}.

\task{Programmet}
\begin{asmcode}
main:
    sbi DDRB, 0                     ; PB0 is an output.
    sbi PORTB, 0                    ; The pulse starts.

delay_start:
    ldi r18, 200
delay_loop:
    dec r18
    brne delay_loop
delay_end:
    cbi PORTB, 0                    ; The pulse ends.

end:
    rjmp end
\end{asmcode}

Programmet ger en puls på PB0: stiftet går högt, fördröjningen körs, och stiftet går lågt.

\task{Uppgifter}
\begin{enumerate}
  \item Räkna ut fördröjningens längd, från \code{delay_start} till \code{delay_end}, i
    klockcykler och i mikrosekunder. Använd en tabell med en rad per instruktion, som i
    \secref{c9:b4}.
  \item Mät samma sak i simulatorn: brytpunkt på \code{delay_start} och \code{delay_end},
    Frequency på 16~MHz, och stoppuret nollställt vid den första brytpunkten. Stämmer det?
  \item Pulsen ska vara \textbf{30~µs}. Räkna ut hur många varv loopen ska gå, ändra programmet, och
    mät.
  \item Vilken är den längsta fördröjning loopen kan ge? Räkna först, mät sedan.
  \item Flytta den första brytpunkten till \code{sbi PORTB, 0}. \textbf{Förutsäg} hur mycket
    mätningen ändras.
\end{enumerate}

\task{Frågor}
\begin{itemize}
  \item Varför blir det exakt $3 \cdot n$ cykler och inte $3 \cdot n + 1$?
\end{itemize}

\redovisa{uträkningen och mätningen för 30~µs.}

\labtask{Del 16}{Längre tidsfördröjning}\phantomsection\label{lab3:d16}
\task{Mål}
Göra en längre fördröjning med två loopar i varandra, och träffa en önskad tid exakt.

\task{Bakgrund}
\Secref{c9:b5}.

\task{Specifikation}
PB0 ska byta läge var \textbf{10:e millisekund}: hög i 10~ms, låg i 10~ms, och så vidare, för
alltid. Byt läge genom att läsa \code{PORTB}, invertera bit~0 med \code{eor} och en mask, och
skriva tillbaka. Fördröjningen mellan bytena görs med två loopar i varandra.

\task{Uppgifter}
\begin{enumerate}
  \item Rita flödesplanen: en loop som byter läge på PB0 och sedan väntar, med fördröjningen som
    två loopar i varandra.
  \item Välj OUTER och INNER med formeln $\text{OUTER} \cdot (3 \cdot \text{INNER} + 3)$ så att
    fördröjningen blir så nära 10~ms som möjligt. Skriv programmet och mät fördröjningen.
  \item Lägg till en \code{nop} i den yttre loopen. Hur ändras formeln? Välj nya konstanter så att
    fördröjningen blir \textbf{exakt} 10~ms, 160\,000 cykler. Mät.
  \item Mät nu en hel halvperiod: från ett byte av PB0 till nästa. Den blir några cykler längre än
    fördröjningen. Hur många, och vilka instruktioner är det?
  \item Kör programmet fritt med \textbf{F5} i några sekunder och stoppa. Hur lång tid har gått
    enligt stoppuret? Ungefär hur många gånger har PB0 bytt läge?
\end{enumerate}

\task{Frågor}
\begin{itemize}
  \item Skulle en lysdiod på PB0 synas blinka med den här fördröjningen på ett riktigt kort?
    Varför inte?
\end{itemize}

\redovisa{konstanterna för exakt 10~ms och mätningarna.}

\labtask{Del 17}{Flervalssituationer och omarbetning av flödesplaner}%
\phantomsection\label{lab3:d17}
\task{Mål}
Skriva ett program som väljer mellan flera alternativ, och arbeta om en flödesplan så att
programmet får en väg in och en väg ut.

\task{Bakgrund}
\Secref{c9:a7} och~\ref{c9:a8}.

\task{Specifikation}
Ett läge i \code{r16} väljer ett mönster på port B:

\begin{narrowtable}{@{}cl@{}}
  \thead{Läge} & \thead{Mönster på port B} \\
  \midrule
  0 & alla släckta \\
  1 & alla tända \\
  2 & de fyra nedre, PB0--PB3 \\
  3 & de fyra övre, PB4--PB7 \\
  annat & varannan, \code{0101 0101}, som felindikering \\
\end{narrowtable}

Programmet går i en loop, så att läget kan ändras i Processor Status medan programmet står
stilla.

\task{Uppgifter}
\begin{enumerate}
  \item Rita en \textbf{första} flödesplan, hur ni vill. Skriv programmet efter den, och testa alla
    lägen, inklusive ett felläge.
  \item Räkna i er flödesplan och i programmet: hur många rutor eller instruktioner skriver till
    port B? Hur många ställen går programmet tillbaka till loopens början från?
  \item \textbf{Arbeta om} flödesplanen enligt \secref{c9:a8}: låt varje alternativ bara välja
    mönster, och låt alla alternativ mötas i \textbf{en} ruta som skriver till port B.
  \item Skriv om programmet efter den nya flödesplanen. Testa alla lägen igen. Programmet ska göra
    exakt samma sak som förut.
  \item Lägg till läge~4, PB0 och PB7 tända. Hur många ställen i programmet behövde ändras, i den
    första versionen och i den omarbetade?
\end{enumerate}

\task{Frågor}
\begin{itemize}
  \item Varför är det viktigt att ha med alternativet \enquote{annat} i flödesplanen?
\end{itemize}

\redovisa{båda flödesplanerna och det omarbetade programmet.}

% Labb 3, parts 18 to 23, from labs/lab3/d_io_and_subroutines.md. Each part is a \labtask, and its
% Mål, Bakgrund, Programmet, Uppgifter and Frågor are \task headings. The programs given in parts
% 20 (with its two constants left to fill in) and 23 (with its deliberate fault) are the
% handout's. One deliberate difference: the last comment of the program in part 23 says to copy
% delay_ms from part 22, where the guide says L10 Appendix C.4, a name the book does not use.

\section{Del 18--23: Text, 16 bitar, inporten och subrutiner}
\label{lab3:sec:d}
Delarna hör till kapitel~\ref{c10:ch}. Läs \secref{c10:app:a} före del 18--20, \secref{c10:app:b}
före del~21 och \secref{c10:app:c} före del 22--23. Instruktionerna finns på referensbladet i
bilaga~\ref{app:avr}.

Del 21 krävs för godkänt. Del 18--20 och 22--23 krävs för VG, och del~22 behövs dessutom för
slutuppgiften i kapitel~\ref{c11:ch}, så gör den även om ni inte siktar på VG.

\labtask{Del 18}{ASCII-koder för hexadecimala siffertecken}\phantomsection\label{lab3:d18}
\task{Mål}
Göra om värden till de tecken som skriver ut dem, och tillbaka, och lägga text i minnet.

\task{Bakgrund}
ASCII-tabellen och luckan på sju mellan \code{'9'} och \code{'A'} finns i
\secref{c10:a1}--\ref{c10:a2}. En byte blir två tecken med \code{swap} och \code{andi} enligt
\secref{c10:a3}, och \code{sts} beskrivs i \secref{c10:a5}. Kom ihåg att \code{subi r17, -0x30}
adderar \code{0x30}.

\task{Uppgifter}
\begin{enumerate}
  \item \textbf{En siffra blir ett tecken.} Ladda \code{r17} med 7 och gör om det till tecknet
    \code{'7'} med en enda instruktion. \emph{Förutsäg} \code{r17} hexadecimalt innan ni stegar.
  \item \textbf{En hexadecimal siffra blir ett tecken.} Ladda \code{r18} med ett värde 0--15,
    kopiera det till \code{r19} och gör om \code{r19} till tecknet för siffran, med \code{cpi},
    \code{brlo} och två \code{subi}. Fyll i tabellen \textbf{innan} ni kör programmet, och
    kontrollera sedan varje rad genom att byta värdet i \code{ldi}:

    \begin{filltable}{@{}p{5em}YY@{}}
      \thead{\code{r18}} & \thead{\code{r19}, förutsagt} & \thead{\code{r19}, i simulatorn} \\
      \midrule
      \code{0x00} & & \fillrule
      \code{0x09} & & \fillrule
      \code{0x0A} & & \fillrule
      \code{0x0B} & & \fillrule
      \code{0x0F} & & \\
    \end{filltable}
  \item \textbf{En byte blir text.} Ladda \code{r16} med \code{0xA7}. Gör om den höga nibblen till
    ett tecken i \code{r20} och den låga till ett tecken i \code{r21}, och lägg dem på adress
    \code{0x0100} och \code{0x0101} med \code{sts}. Öppna fönstret \textbf{Memory}, välj
    \textbf{data IRAM} och gå till adress \code{0x0100}. Står texten \code{A7} i ASCII-kolumnen
    till höger? Prova också \code{0x3C}, \code{0xFF} och \code{0x00}.
  \item \textbf{Tillbaka från tecken till värde.} Ladda \code{r22} med \code{'F'} och gör om det
    till värdet 15, enligt \secref{c10:a4}. Prova också \code{'0'}, \code{'9'} och \code{'A'}.
\end{enumerate}

\task{Frågor}
\begin{itemize}
  \item Varför måste den höga nibblen flyttas ned med \code{swap} innan den kan göras om till ett
    tecken?
  \item Vad blir tecknet om man glömmer \code{andi r20, 0x0F} efter \code{swap}, för byten
    \code{0xA7}?
\end{itemize}

\redovisa{uppgift 3 och 4 i simulatorn, med texten synlig i Memory-fönstret.}

\labtask{Del 19}{16-bitarsoperationer}\phantomsection\label{lab3:d19}
\task{Mål}
Räkna med tal som inte får plats i en byte: addera, subtrahera, räkna och jämföra 16-bitarstal.

\task{Bakgrund}
Registerpar och \code{low()}/\code{high()} finns i \secref{c10:a6}, addition och subtraktion med
minnessiffra i \secref{c10:a7}, \code{adiw}, \code{sbiw} och \code{movw} i \secref{c10:a8} och
jämförelsen med \code{cpi}/\code{cpc} i \secref{c10:a9}.

\task{Uppgifter}
\begin{enumerate}
  \item \textbf{Addition.} Beräkna 1000 + 2000 med \code{r25:r24} och \code{r23:r22}, med
    \code{add} och \code{adc}. Kopiera resultatet till \code{r19:r18} med \code{movw}.

    \emph{Förutsäg} först, med additionen uppställd för hand: vad blir \code{r24} och flaggan
    \code{C} efter \code{add}, och vad blir \code{r25} efter \code{adc}? Stega sedan och jämför i
    Processor Status.
  \item \textbf{Subtraktion.} Beräkna 3000 -- 1234 med \code{sub} och \code{sbc}, och lägg
    resultatet i \code{r21:r20}. \emph{Förutsäg} båda bytena hexadecimalt. Blir det ett lån från
    den låga byten?
  \item \textbf{En 16-bitars räknare.} Nollställ \code{r25:r24} och räkna upp det med \code{adiw}
    tills det är 1000, med \code{cpi} och \code{cpc} i loopen. Sätt en etikett direkt efter loopen.

    \emph{Förutsäg} hur många cykler loopen tar, med samma metod som i kapitel~\ref{c9:ch}. Mät
    sedan med brytpunkter före och efter loopen och fönstret Processor Status
    (\secref{studio:4-4}). Stämmer siffrorna? Om inte: vilket varv räknade ni fel på?
  \item \textbf{När 16 bitar inte räcker.} Ladda \code{r27:r26} med \code{0xFFFF} och addera 1 med
    \code{adiw}. \emph{Förutsäg} resultatet och flaggorna \code{C} och \code{Z}.
\end{enumerate}

\task{Frågor}
\begin{itemize}
  \item Varför måste de låga bytena adderas först?
  \item Varför behövs registret med \code{high(1000)} i uppgift~3, när den låga byten kan jämföras
    direkt med \code{cpi}?
\end{itemize}

\redovisa{uppgift 1--4, och er förutsägelse och mätning i uppgift~3.}

\labtask{Del 20}{Lång tidsfördröjning}\phantomsection\label{lab3:d20}
\task{Mål}
Konstruera en fördröjning på en halv sekund, exakt på cykeln, och få en lysdiod att blinka en gång
per sekund.

\task{Bakgrund}
Metoden, en 16-bitars inre loop i en yttre 8-bitars loop, och formeln \code{M(4N + 4)} finns i
\secref{c10:a10}. Exemplet där blinkar med en sekunds halvperiod; här ska det gå dubbelt så fort.

\task{Programmet}
Fyll i de två konstanterna.

\begin{asmcode}
.include "m328Pdef.inc"

.equ OUTER_TURNS = ?                ; Work these two out before you run anything.
.equ INNER_TURNS = ?

.org 0x0000
    rjmp main

main:
    sbi DDRB, 5                     ; PB5, Arduino pin 13, is an output.
    ldi r17, 0b00100000             ; The mask that selects bit 5.

loop:
    in r16, PORTB                   ; Toggle PB5.
    eor r16, r17
    out PORTB, r16

delay_begin:
    ldi r20, OUTER_TURNS
delay_outer:
    ldi r24, low(INNER_TURNS)
    ldi r25, high(INNER_TURNS)
delay_inner:
    sbiw r24, 1
    brne delay_inner
    dec r20
    brne delay_outer
delay_end:
    rjmp loop
\end{asmcode}

\task{Uppgifter}
\begin{enumerate}
  \item \textbf{Räkna först.} Hur många cykler är en halv sekund vid 16~MHz? Välj
    \code{INNER_TURNS} så att ett varv i den yttre loopen tar 200\,000 cykler, och räkna ut
    \code{OUTER_TURNS}. Skriv upp uträkningen.
  \item \textbf{Mät.} Sätt stoppurets frekvens till 16~MHz, sätt en brytpunkt på
    \code{delay_begin} och en på \code{delay_end}, och mät tiden mellan dem. \emph{Förutsäg}
    stoppurets värde innan ni kör.
  \item \textbf{Hela perioden.} Hur lång är en hel blinkperiod, från en tändning till nästa, räknat
    i cykler? Varför är den inte exakt en sekund? Hur mycket fel är den, i mikrosekunder?
  \item \textbf{En annan tid.} Ändra till en fjärdedels sekund. Vilken av de två konstanterna ändrar
    ni, och till vad?
\end{enumerate}

\begin{aside}
\textbf{Tips.} Simulatorn behöver en stund för att köra åtta miljoner cykler. Använd
\textbf{Continue} (\textbf{F5}) mellan brytpunkterna, inte \textbf{Step Into}.
\end{aside}

\task{Frågor}
\begin{itemize}
  \item Varför är det enklare att välja ett \enquote{jämnt} antal cykler för den yttre loopen först,
    och räkna ut antalet varv sedan?
  \item Vad händer om \code{OUTER_TURNS} sätts till 0?
\end{itemize}

\redovisa{uträkningen och mätningen.}

\labtask{Del 21}{Inporten PORTD (PIND)}\phantomsection\label{lab3:d21}
\task{Mål}
Läsa knappar på port D och styra lysdioder på port B utifrån dem, och bygga logiska funktioner i
mjukvara.

\task{Bakgrund}
Pull-up-motståndet och varför en nedtryckt knapp läses som 0 finns i \secref{c10:b2}, att läsa
hela porten i \secref{c10:b3} och \code{sbic}/\code{sbis} i \secref{c10:b4}. Så simuleras en
knapp: \secref{studio:4-5}.

I hela delen sitter knapparna mellan stiftet och jord, och de inbyggda pull-uparna används. En
\textbf{tom} ruta i \code{PIND} är alltså en \textbf{nedtryckt} knapp.

\task{Uppgifter}
\begin{enumerate}
  \item \textbf{Knapparna på lysdioderna.} Skriv av programmet \code{buttons_to_leds.asm} från
    \secref{c10:b3}. Kör till loopen, stoppa, och ändra \code{PIND} för hand. \emph{Förutsäg}
    \code{PORTB} för varje rad innan ni stegar:

    \begin{filltable}{@{}p{6.4em}YYY@{}}
      \thead{\code{PIND}} & \thead{Nedtryckta knappar} & \thead{\code{PORTB}, förutsagt} &
        \thead{\code{PORTB}, i simulatorn} \\
      \midrule
      \code{0b11111111} & & & \fillrule
      \code{0b11111011} & & & \fillrule
      \code{0b11110011} & & & \fillrule
      \code{0b01111110} & & & \\
    \end{filltable}
  \item \textbf{En knapp, en lysdiod.} Skriv ett program där lysdioden på PB5 lyser så länge
    knappen på PD2 hålls nedtryckt, med \code{sbic} eller \code{sbis} och utan att läsa hela porten.
  \item \textbf{AND och OR i mjukvara.} Två knappar, på PD2 och PD3. Lysdioden på PB0 ska lysa när
    \textbf{båda} är nedtryckta, och lysdioden på PB1 när \textbf{minst en} är det. Det är samma två
    funktioner som seriekopplingen och parallellkopplingen i Labb~1 (bilaga~\ref{app:lab1}).

    Rita en flödesplan först. Fyll sedan i sanningstabellen, med förutsägelsen före simuleringen:

    \begin{filltable}{@{}p{5.6em}p{5.6em}YY@{}}
      \thead{PD2} & \thead{PD3} & \thead{PB0 (AND)} & \thead{PB1 (OR)} \\
      \midrule
      släppt & släppt & & \fillrule
      släppt & nedtryckt & & \fillrule
      nedtryckt & släppt & & \fillrule
      nedtryckt & nedtryckt & & \\
    \end{filltable}
  \item \textbf{Extra.} Lägg till lysdioden på PB2, som ska lysa när \textbf{exakt en} av knapparna
    är nedtryckt. Vilken grind från kapitel~\ref{c2:ch} är det? Vilken koppling i Labb~1 gjorde
    samma sak?
\end{enumerate}

\task{Frågor}
\begin{itemize}
  \item Vad läser programmet på en ingång utan pull-up, när knappen är släppt?
  \item I uppgift~3 behövs \code{com} och \code{andi} innan jämförelsen. Vad gör var och en av dem,
    och vad går fel om man hoppar över \code{andi}?
  \item Varför är det enklare att ändra funktionen i uppgift~3 än i Labb~1?
\end{itemize}

\redovisa{uppgift~3, med sanningstabellen, i simulatorn.}

\labtask{Del 22}{Subrutiner}\phantomsection\label{lab3:d22}
\task{Mål}
Flytta en fördröjning till en subrutin, se returadressen på stacken, och skriva en subrutin som tar
en parameter.

\task{Bakgrund}
\code{rcall}, \code{ret}, stacken och stackpekaren finns i \secref{c10:c2}--\ref{c10:c3},
subrutinen \code{delay_ms} med parameter i \secref{c10:c4}, och en subrutin som anropar en annan i
\secref{c10:c7}.

Från och med den här delen ska varje program som använder \code{rcall} börja med att initiera
stackpekaren, med de fyra raderna i \secref{c10:c3}.

\task{Uppgifter}
\begin{enumerate}
  \item \textbf{Fördröjningen blir en subrutin.} Ta programmet från del~20. Flytta fördröjningen,
    från \code{delay_begin} till och med \code{brne delay_outer}, till en subrutin
    \code{delay_500ms} som slutar med \code{ret}, och anropa den med \code{rcall} där fördröjningen
    stod. Lägg till initieringen av \code{SP}. Programmet ska blinka precis som förut.
  \item \textbf{Returadressen.} \emph{Förutsäg} \code{SP} före \code{rcall}, i början av
    subrutinen och efter \code{ret}. Leta upp adressen till instruktionen efter \code{rcall} i
    listfilen (\code{.lss}, \secref{c7:b7}). Stega in i subrutinen med \textbf{F11} och läs de två
    översta bytena i stacken i fönstret Memory (\textbf{data IRAM}, adress
    \code{0x08FE}--\code{0x08FF}). Är det returadressen? Vilken byte ligger var?
  \item \textbf{En parameter.} Skriv av \code{delay_ms} från \secref{c10:c4}. Skriv om
    \code{delay_500ms} så att den laddar \code{r24} med 250 och anropar \code{delay_ms} två gånger.
    Glöm inte att \code{delay_500ms} själv ändrar \code{r24}, och alltså ska spara det.
  \item \textbf{Järnvägsövergången.} Ändra huvudprogrammet så att lysdioderna på PB0 och PB1
    blinkar växelvis, en halv sekund var, som vid en järnvägsövergång.
  \item \textbf{Mät ett anrop.} \emph{Förutsäg} hur många cykler ett anrop av \code{delay_500ms}
    tar, från \code{rcall} till instruktionen efter, med formeln för \code{delay_ms} och cyklerna
    för \code{rcall}, \code{push}, \code{ldi}, \code{pop} och \code{ret}. Mät sedan med
    stoppuret.
  \item \textbf{Den djupaste punkten.} Stoppa programmet mitt i \code{delay_ms}, när den har
    anropats från \code{delay_500ms}. \emph{Förutsäg} \code{SP} först, och kontrollera. Hur många
    byte ligger på stacken, och vad är var och en?
\end{enumerate}

\task{Frågor}
\begin{itemize}
  \item Vad skulle hända om \code{delay_ms} inte sparade \code{r24}?
  \item \code{rcall} kostar 3 cykler och \code{ret} 4. Är det ett argument mot att använda
    subrutiner?
\end{itemize}

\redovisa{uppgift 2, 4 och 5.}

\labtask{Del 23}{Använda stacken som lagringsplats}\phantomsection\label{lab3:d23}
\task{Mål}
Använda \code{push} och \code{pop} för att spara register, förstå varför ordningen spelar roll, och
se vad som händer när stacken inte stämmer.

\task{Bakgrund}
\code{push}, \code{pop} och sist in, först ut finns i \secref{c10:c5}, och kursens regel om att
spara det man ändrar i \secref{c10:c6}.

\task{Programmet}
Huvudprogrammet räknar i \code{r16} hur många gånger det har anropat \code{flash_all}, och slutar
efter tre gånger. Subrutinen \code{flash_all} tänder alla lysdioder på port B i 100~ms och släcker
dem i 100~ms, och använder själv \code{r16} till något helt annat. Den saknar \code{push} och
\code{pop}, med flit.

\begin{asmcode}
.include "m328Pdef.inc"

.org 0x0000
    rjmp main

main:
    ldi r16, high(RAMEND)
    out SPH, r16
    ldi r16, low(RAMEND)
    out SPL, r16

    ldi r16, 0xFF
    out DDRB, r16                   ; Port B: all outputs.

    clr r16                         ; r16 counts the flashes.
loop:
    inc r16
    rcall flash_all
    cpi r16, 3
    brne loop
    out PORTB, r16                  ; Show the count.
end:
    rjmp end

flash_all:
    ldi r16, 0xFF                   ; All LEDs on ...
    out PORTB, r16
    ldi r24, 100
    rcall delay_ms                  ; ... for 100 ms ...
    clr r16                         ; ... and off ...
    out PORTB, r16
    ldi r24, 100
    rcall delay_ms                  ; ... for 100 ms.
    ret

; Copy delay_ms from part 22 here.
\end{asmcode}

\task{Uppgifter}
\begin{enumerate}
  \item \textbf{Byt plats med stacken.} Ladda \code{r20} med \code{0x55} och \code{r21} med
    \code{0xAA} i ett eget litet program. Lägg dem på stacken med \code{push r20}, \code{push r21},
    och hämta dem med \code{pop r20}, \code{pop r21}. \emph{Förutsäg} värdena i \code{r20} och
    \code{r21} efteråt, och \code{SP} efter varje instruktion.
  \item \textbf{Felet.} Kör programmet ovan. \emph{Förutsäg} först vad det gör. Hur många gånger
    blinkar det? Stega igenom ett varv i loopen och följ \code{r16}. Förklara vad som går fel.
  \item \textbf{Rättningen.} Lägg till \code{push} och \code{pop} i \code{flash_all} så att
    subrutinen följer kursens regel. Vilka register måste sparas? Programmet ska nu blinka tre
    gånger och sedan visa talet 3 på lysdioderna.
  \item \textbf{Djupet.} Stoppa programmet mitt i den första \code{delay_ms}. \emph{Förutsäg}
    \code{SP}, och kontrollera. Hur många byte ligger på stacken, och vad är var och en?
  \item \textbf{Fel ordning.} Byt plats på de två \code{pop} i \code{flash_all}, så att de inte
    längre står i omvänd ordning mot \code{push}. Vad händer? Förklara med stacken.
  \item \textbf{En pop för lite.} Ta bort en av \code{pop}-instruktionerna. Stega fram till
    \code{ret} i \code{flash_all} och läs \code{SP}. Stega en gång till: vart hoppar programmet?
    Förklara med stacken.
\end{enumerate}

\task{Frågor}
\begin{itemize}
  \item Varför måste \code{pop} göras i omvänd ordning mot \code{push}?
  \item Felet i uppgift~6 syns tydligare än felet i uppgift~5. Varför är det i själva verket en
    fördel?
\end{itemize}

\redovisa{uppgift 2, 3 och 4.}

% Labb 3, the final task, parts A and B, from labs/lab3/e_final_task.md. The wiring and the
% ready-made subroutines, headings of their own in the guide, are subsections; parts A and B are
% \labtask, with their Mål, Bakgrund, Specifikation, Uppgifter, Frågor and Extra as \task
% headings. The subroutines delay_s and delay_ms are the handout's, and are printed as the guide
% prints them.

\section{Slutuppgift: trafikljuset och övergångsstället}
\label{lab3:sec:e}
Slutuppgiften hör till kapitel~\ref{c11:ch}, och är den del av Labb~3 där programmet till sist
lämnar simulatorn och styr något på riktigt: tre lysdioder och en knapp på ett Arduino Uno-kort.
Del~A krävs för godkänt, och del~B för VG.

Läs \secref{c11:app:a} om lysdioder, knappar och kopplingen, och \secref{c11:app:b}, där ett
enklare ljus, ett vägarbetsljus, tas hela vägen från flödesplan till program. Hur programmet förs
över till kortet står i \secref{studio:5}.

\subsection{Kopplingen}
\label{lab3:sec:wiring}
Samma koppling används i båda delarna. Varje lysdiod har ett eget förkopplingsmotstånd, och knappen
kopplas mellan stift 2 och jord.

\begin{narrowtable}{@{}lcc@{}}
  \thead{Funktion} & \thead{Arduino-stift} & \thead{Port och bit} \\
  \midrule
  Röd lysdiod & 8 & PB0 \\
  Gul lysdiod & 9 & PB1 \\
  Grön lysdiod & 10 & PB2 \\
  Knapp, till jord & 2 & PD2 \\
  Jord & GND & -- \\
\end{narrowtable}

\bookfigure{l11_traffic_wiring}{Kopplingsschemat: tre lysdioder med förkopplingsmotstånd på stift
8, 9 och 10, och en knapp från stift 2 till jord.}{lab3:fig:wiring}

Lysdiodens långa ben, anoden, ska vara vänt mot motståndet och stiftet, och det korta, katoden, mot
jord. Koppla med USB-kabeln urkopplad, och kontrollera kopplingen mot tabellen innan kabeln sätts
i.

\subsection{De färdiga subrutinerna}
\label{lab3:sec:subroutines}
Båda delarna använder fördröjningar på hela sekunder. Ni behöver inte skriva dem själva: kopiera in
de två subrutinerna nedan sist i programmet, och konstanten \code{QUARTER_MS} överst. De förklaras
i \secref{c11:b4}, och \code{delay_ms} är samma subrutin som i del~22.

\begin{asmcode}
.equ QUARTER_MS = 250               ; 250 ms. Set it to 1 to run about 250 times
                                    ; faster in the simulator; set it back before
                                    ; flashing.
\end{asmcode}

\begin{asmcode}
;--------------------------------------------------------------------------------
; delay_s: Waits r24 seconds (1-63), as 4 x r24 calls of delay_ms with 250 ms each.
; A call takes 16 000 084 x r24 + 18 cycles: about 5 microseconds too long per second.
;--------------------------------------------------------------------------------
delay_s:
    push r24                        ; Save the registers this subroutine changes.
    push r25
    mov r25, r24
    lsl r25                         ; r25 = 2 x r24 ...
    lsl r25                         ; ... = 4 x r24: the number of quarter seconds.
    ldi r24, QUARTER_MS             ; A quarter of a second is 250 ms.
delay_s_loop:
    rcall delay_ms
    dec r25
    brne delay_s_loop
    pop r25                         ; Restore them, in the opposite order.
    pop r24
    ret

;--------------------------------------------------------------------------------
; delay_ms: Waits r24 milliseconds (1-255) at 16 MHz.
; A call takes 16 000 x r24 + 18 cycles, the rcall and the ret included.
;--------------------------------------------------------------------------------
delay_ms:
    push r24                        ; Save the registers this subroutine changes.
    push r26
    push r27
delay_ms_outer:
    ldi r26, low(3999)              ; X = r27:r26 = 3999.
    ldi r27, high(3999)
delay_ms_inner:
    sbiw r26, 1                     ; 4 cycles per turn, 3 the last time.
    brne delay_ms_inner
    dec r24                         ; One more millisecond done.
    brne delay_ms_outer
    pop r27                         ; Restore them, in the opposite order.
    pop r26
    pop r24
    ret
\end{asmcode}

Så används de: ladda \code{r24} med antalet sekunder och anropa \code{delay_s}. \code{ldi r24, 3}
följt av \code{rcall delay_s} väntar tre sekunder, och \code{r24} är fortfarande 3 efteråt.

\begin{aside}
\textbf{I simulatorn} tar varje simulerad sekund lång tid att köra. Sätt \code{QUARTER_MS} till 1
medan ni testar programmet där: då går varje \enquote{sekund} på ungefär 4~ms, och hela ljuscykeln
syns på några ögonblick. \textbf{Sätt tillbaka den till 250 innan programmet förs över till
kortet.}
\end{aside}

\labtask{Del A}{Trafikljuset}\phantomsection\label{lab3:a}
\task{Mål}
Styra ett trafikljus med tre lysdioder genom en fast ljuscykel, i simulatorn och på kortet.

\task{Bakgrund}
Vägarbetsljuset i \secref{c11:app:b} har två tillstånd; trafikljuset har fyra. Metoden är
densamma: skriv upp tillstånden och deras tider, rita flödesplanen, och skriv sedan programmet ett
tillstånd i taget.

\task{Specifikation}
Trafikljuset går igenom fyra tillstånd, i den här ordningen, och börjar sedan om:

\begin{narrowtable}{@{}lcccc@{}}
  \thead{Tillstånd} & \thead{Röd (PB0)} & \thead{Gul (PB1)} & \thead{Grön (PB2)} & \thead{Tid} \\
  \midrule
  Stopp & tänd & släckt & släckt & 3~s \\
  Klart att köra & tänd & tänd & släckt & 1~s \\
  Kör & släckt & släckt & tänd & 3~s \\
  Stopp strax & släckt & tänd & släckt & 1~s \\
\end{narrowtable}

\bookfigure{l11_traffic_timing}{Tidsdiagram för trafikljuset: en ljuscykel är 8~s.}%
{lab3:fig:timing}

\task{Uppgifter}
\begin{enumerate}
  \item \textbf{Flödesplan.} Rita flödesplanen för programmet, med ett symbolpar per tillstånd:
    tänd rätt lampor, vänta rätt tid. \kindtag[fillaccent]{Redovisa} den innan ni börjar skriva kod.
  \item \textbf{Bitmönstren.} Skriv upp vilket värde \code{PORTB} ska ha i vart och ett av de fyra
    tillstånden, binärt och hexadecimalt. Använd gärna \code{.equ RED = 0} och skrivsättet
    \code{(1 << RED)}, så att programmet säger vad det gör.
  \item \textbf{Programmet.} Skriv programmet utifrån programmallen, med initiering av \code{SP},
    \code{DDRB}, och de färdiga subrutinerna.
  \item \textbf{Simulera.} Sätt \code{QUARTER_MS} till 1 och kör i simulatorn. Sätt en brytpunkt på
    varje \code{out PORTB, r16} och kontrollera i I/O-fönstret att rätt lampor tänds, i rätt
    ordning. Kontrollera också med stoppuret att tiderna står i rätt förhållande till varandra: 3,
    1, 3, 1.
  \item \textbf{På kortet.} Sätt tillbaka \code{QUARTER_MS} till 250, bygg, och för över programmet
    till kortet enligt \secref{studio:5}.
  \item \textbf{Mät.} Ta tid på tio hela cykler med stoppuret i en mobiltelefon. \emph{Förutsäg}
    tiden först, med formeln för \code{delay_s}. Hur nära kommer ni? Vad begränsar noggrannheten i
    er mätning, och vad begränsar noggrannheten i programmet?
\end{enumerate}

\task{Frågor}
\begin{itemize}
  \item Vad händer om \code{delay_s} inte sparade \code{r24}? Hade programmet märkt det?
  \item Varför tänds röd och gul samtidigt i det andra tillståndet? Vad säger det till en bilist?
\end{itemize}

\redovisa{flödesplanen, simuleringen och trafikljuset på kortet.}

\labtask{Del B}{Övergångsstället}\phantomsection\label{lab3:b}
\task{Mål}
Låta en knapp styra när ljuscykeln körs: ett övergångsställe där bilarna har grönt tills en
fotgängare trycker på knappen.

\task{Bakgrund}
Att vänta på en knapp med \code{sbic} finns i \secref{c10:b5}, och varför en knapp studsar i
\secref{c11:a4}. Subrutiner och kursens regel om register finns i \secref{c10:app:c}.

\task{Specifikation}
\begin{itemize}
  \item Bilarna har \textbf{grönt}, så länge ingen trycker på knappen.
  \item När knappen trycks går ljuset igenom:

    \begin{narrowtable}{@{}lcccc@{}}
      \thead{Tillstånd} & \thead{Röd (PB0)} & \thead{Gul (PB1)} & \thead{Grön (PB2)} &
        \thead{Tid} \\
      \midrule
      Stopp strax & släckt & tänd & släckt & 1~s \\
      Stopp: fotgängarna går & tänd & släckt & släckt & 4~s \\
      Klart att köra & tänd & tänd & släckt & 1~s \\
    \end{narrowtable}
  \item Därefter blir det grönt igen, och programmet väntar på nästa tryck.
  \item Knappen är aktivt låg, med den inbyggda pull-upen: PD2 som ingång, med pull-up.
  \item Programmet ska använda minst en egen subrutin utöver de färdiga, och varje subrutin ska
    följa kursens regel: spara varje register den ändrar.
\end{itemize}

\task{Uppgifter}
\begin{enumerate}
  \item \textbf{Flödesplan.} Rita flödesplanen, med beslutet \enquote{knappen nedtryckt?} och de
    tre tillstånden. \kindtag[fillaccent]{Redovisa} den innan ni skriver kod.
  \item \textbf{En egen subrutin.} De tre tillstånden gör samma sak med olika värden: tänd vissa
    lampor, vänta en viss tid. Skriv en subrutin \code{show} som tänder lamporna i \code{r16} och
    väntar det antal sekunder som står i \code{r24}. Vilka register ändrar den? Måste något av dem
    sparas?
  \item \textbf{Programmet.} Skriv programmet. Initiera \code{SP}, \code{DDRB}, \code{DDRD} och
    pull-upen på PD2.
  \item \textbf{Simulera.} Sätt \code{QUARTER_MS} till 1. Kör programmet, stoppa det medan det
    väntar på knappen, och \enquote{tryck} genom att tömma rutan för bit~2 i \code{PIND}.
    \emph{Förutsäg} vilka lampor som ska tändas i vilken ordning, och kontrollera med
    brytpunkter.
  \item \textbf{På kortet.} Sätt tillbaka \code{QUARTER_MS} till 250 och för över programmet till
    kortet. Prova ett kort tryck, ett långt tryck, och ett tryck medan ljuset redan är rött.
\end{enumerate}

\task{Frågor}
\begin{itemize}
  \item Knappen studsar när den trycks. Varför ställer det inte till något här, när det gjorde det
    för knapptryckräknaren i \secref{c11:a4}?
  \item Vad händer om någon håller knappen nedtryckt hela tiden? Är det ett problem för bilarna?
  \item Ett tryck medan ljuset är rött gör ingenting. Varför inte? Hur skulle programmet behöva
    ändras för att komma ihåg ett sådant tryck?
\end{itemize}

\task{Extra}
Frivilliga utökningar, för den som vill vidare:
\begin{itemize}
  \item \textbf{Minsta tid på grönt.} Riktiga övergångsställen ger bilarna grönt i minst några
    sekunder innan ett nytt tryck får verkan. Låt bilarna alltid ha grönt i minst 5~s efter varje
    fotgängarfas.
  \item \textbf{Fotgängarnas ljus.} Lägg till en röd och en grön lysdiod för fotgängarna på PB3 och
    PB4 (stift 11 och 12). Fotgängarnas gröna ska lysa bara under de 4 sekunderna med rött för
    bilarna.
  \item \textbf{Vänta-lampan.} Tänd en lampa, till exempel PB5 (stift 13, lysdioden på kortet), när
    knappen har tryckts, och släck den när fotgängarna får grönt.
\end{itemize}

\redovisa{flödesplanen, simuleringen och övergångsstället på kortet.}

